\section{Wilson's theorem}
\phantomsection
\label[appendix]{app:wilsons-theorem}

\begin{theorem}[Wilson's theorem] \index{Wilson's theorem}
$p$ is prime if and only if the factorial $(p - 1)! \equiv -1 \pmod p$.
\end{theorem}

where the factorial $(p - 1)! = 1 \times 2 \times 3 \dotsm (p - 1)$. In other words, this factorial plus 1 can be divided by $p$. We first prove a lemma.

\begin{lemma}
For positive integer $a$ less than prime $p$, there exists unique positive integer $b$ less than $p$, such that $ab \equiv 1 \pmod p$. We call $b$ the inverse of $a$ under the multiplication modular $p$.
\end{lemma}

\begin{proof}
$a$ and $p$ are relatively prime, that is, $(a, p) = 1$. By Bézout's identity, there are $x, y$, such that $ax + py = 1$. It follows that $p$ divides $ax - 1$, that is, $ax \equiv 1 \pmod p$. Let $b = x + pk$; adjust the integer $k$ such that $0 < b < p$. Next, we are going to show the uniqueness. Suppose there is some positive integer $b' < p$ satisfying $ab' \equiv 1 \pmod p$, taking their difference:

\[
ab - ab' \equiv a(b - b') \equiv 1 - 1 \equiv 0 \pmod p
\]

Therefore, $p$ divides $a(b - b')$. Since $p$ does not divide $a$, it follows that $p$ divides $b - b'$, but both positive integers $b, b'$ are less than $p$, therefore, $b = b'$.
\end{proof}

We are going to prove Wilson's theorem with this lemma. The idea is to rearrange the factors of the factorial $(p - 1)!$, pair everyone with its inverse under the multiplication modular $p$:

\begin{proof}
By the lemma of inverse modular $p$, for every positive integer $a$ less than $p$, there is unique inverse, namely the positive integer $b$, such that $ab \equiv 1 \pmod p$. Rearrange the factorial:

\begin{align*}
(p - 1)! &= 1 \times 2 \times 3 \dotsm (p - 1)  \\
  &= 1 (a_1 b_1) (a_2 b_2) \dotsm (a_k b_k) (p - 1) && p - 1\text{ is even} \\
  &\equiv 1 \cdot 1 \cdot 1 \dotsm 1 \cdot (p - 1) \pmod p && \text{each } a_ib_i \equiv 1 \pmod p \\
  &\equiv -1 \pmod p
\end{align*}

Conversely, if $p$ is composite, there are three cases: (1) $p = 4$, then $(p - 1)! = 3! = 6$, the remainder is 2 when divided by 4. (2) If $p = q^2$ is a square of prime number, when $q \geq 3$:

\begin{align*}
(q - 1)^2 &\geq (3 - 1)^2 = 4 > 2\\
q^2 - 2q + 1 &> 2 \\
2q &< q^2 - 1 = p - 1
\end{align*}
It follows that both $q, 2q$ are factors of the factorial $(p - 1)!$, that is $2q^2 = 2p$ is a factor, which implies $p$ divides this factorial, but not remains $-1$. (3) $p = ab$, where $a, b$ both greater than 1 and less than $p$, hence they are factors of the factorial $(p - 1)!$. Therefore, $p = ab$ divides the factorial, but does not remain $-1$. Summarizing the three cases, $p$ cannot be composite, but prime.
\end{proof}
